\subsection*{Notes for II. PROBLEM DEFINITION}
\paragraph{Objective}
Define the problem setting at a high level: what is being optimized, under which constraints, and what counts as budget/cost.

\paragraph{Keep}
\begin{itemize}
  \item The focus on evaluation cost and realistic constraints.
\end{itemize}

\textbf{TODO}:
\begin{itemize}
    \item \paragraph{Rewrite / move}
\end{itemize}
\begin{itemize}
  \item \begin{itemize}
      \item Centralize the definitions that currently appear later in your protocol: define the evaluation unit, budgets, and objective measurement rules here (with full math in the formulation section).
  \end{itemize}
\end{itemize}

\begin{itemize}
    \item \paragraph{Add first}
\end{itemize}
\begin{itemize}
      \begin{itemize}
          \item Define what an “evaluation” is (architecture instantiation + training recipe + returned metrics).
          \end{itemize}
          \begin{itemize}
          \item Define budget accounting: whether $B$ counts (i) complete evaluations, (ii) total epochs across ASHA rungs, or (iii) another compute-normalized unit (be explicit).
          \item Define multi-fidelity setting at a conceptual level (ASHA/Hyperband-style partial training).
          \item Define the goal: maximize performance under budget, optionally with multiobjective trade-off.
      \end{itemize}
  \end{itemize}


\paragraph{Peer-review must-haves}
\begin{itemize}
  \item A definition that makes baseline comparisons fair and auditable.
  \item A clear statement of what is held constant across methods (evaluator, search space, constraints).
\end{itemize}
